% Sección 1, Juan Diego (#50). Resumen ejecutivo de una página. Todas las cifras están en las
% secciones que se citan: no hay aquí ninguna cifra nueva.
\section{Resumen ejecutivo}\label{sec:resumen}

\textbf{Problema.} Este trabajo estudia si la historia del hurto a personas en Bogotá permite anticipar cuántas víctimas registrará la Policía Nacional al día siguiente, con menos error que una regla ingenua. Lo que se pronostica es el \emph{registro}: el hurto que no se denuncia ni llega a la autoridad no está en la serie.

\textbf{Datos.} La fuente principal son los 741.986 registros de hurto a personas en Bogotá que la Policía publica a nivel de registro para 2020--2025 (sección~\ref{sec:datos}). Agregados por día forman una serie de 2.131 días, del 2020-01-01 al 2025-10-31, con 730.870 víctimas (sección~\ref{sec:objetivo}). Noviembre y diciembre de 2025 se excluyen porque el registro cae a la mitad sin que lo haga la Línea 123. Las llamadas de la Línea 123 y los procesos de la Fiscalía se usan para comparar tendencias y describir el territorio, sin sumarse a la serie.

\textbf{Método.} Los datos hasta el 2025-04-30 se usan para entrenar, y los 184 días siguientes, de mayo a octubre de 2025, para la evaluación final, que se calcula una sola vez. Las configuraciones se eligen con seis pliegues de origen rodante dentro del entrenamiento. La barra a superar es el ingenuo estacional, que repite el valor de hace una semana, con un MAE de 49,04 víctimas por día (sección~\ref{sec:barra}). Se comparan cinco modelos: SARIMA y ETS, que modelan la propia serie, y una regresión lineal, un bosque aleatorio y una potenciación del gradiente, que usan rezagos y calendario. La prueba de Diebold-Mariano indica si las diferencias superan el azar de una sola ventana.

\textbf{Resultado principal.} SARIMA, ETS y los dos árboles bajan el MAE de la barra a entre 32,72 y 34,22 víctimas por día, entre un 30 y un 33\,\% menos, con $p < 10^{-7}$; la regresión lineal lo baja a 36,31, un 26\,\% menos (tabla~\ref{tab:resultados}). Eso equivale a un error de alrededor del 10\,\% del valor diario. La mejora se repite en los seis pliegues y no depende de cómo se trate 2020. En cambio, entre los cuatro mejores no hay un ganador: los valores $p$ van de 0,17 a 0,91, y la regresión solo se distingue del gradiente si no se corrige la autocorrelación de los errores (sección~\ref{sec:res-entre-modelos}). Los modelos se apoyan en lo mismo que mostró el análisis exploratorio: el ciclo semanal, con un viernes un 65\,\% más alto que el domingo, el efecto de los festivos y el nivel del día anterior (secciones~\ref{sec:eda-semana} y~\ref{sec:interpretacion}).

\textbf{Limitación principal.} La evaluación final no incluye la temporada de noviembre a enero, cuando el ciclo semanal se rompe. En los pliegues que la cubren, el MAE de los modelos llega a 44,61, así que para esas semanas cabe esperar errores mayores. Además, el pronóstico es de toda la ciudad, no trae intervalos de predicción y hereda cualquier cambio en la forma de registrar. Y no está pensado para asignar vigilancia por zonas (sección~\ref{sec:limitaciones}).
